\begin{lemma}[Power of a power in a monoid for nonnegative integer exponents] \label{lemma:power_of_power_in_monoid_for_nonnegative_integers}
    \TODO{AI generated, verify}
    Let $(M, \cdot)$ be a \CrefAndHyperrefIfExist{definition:monoid}{monoid} with identity element $e$, and let $a \in M$. For all nonnegative integers $m, n \in \bbZ_{\geq 0}$,
    $$ (a^m)^n = a^{mn}, $$
    where $a^k$ denotes the \CrefAndHyperrefIfExist{definition:power_of_a_monoid_element_by_a_nonnegative_integer}{$k$th power of $a$}.
\end{lemma}

\begin{proof}
    \TODO{AI generated, verify}
    We proceed by case analysis on whether $m = 0$.

    If $m = 0$, then $a^m = a^0 = e$ by \CrefIfExists{definition:power_of_a_monoid_element_by_a_nonnegative_integer}. Hence $(a^m)^n = e^n = e$ and $a^{mn} = a^0 = e$, so the identity holds.

    If $m > 0$, then for any nonnegative integer $k$, $a^k$ coincides with the power of $a$ in the underlying semigroup\CrefIfExists{definition:power_of_a_monoid_element_by_a_nonnegative_integer}. Thus $(a^m)^n = a^{mn}$ follows from \Cref{lemma:power_of_power_in_semigroup_for_positive_integers} when $n > 0$, and when $n = 0$ both sides equal $e$.
\end{proof}
